\documentclass[11pt]{article}

\usepackage[T1]{fontenc}
\usepackage[utf8]{inputenc}
\usepackage{lmodern}
\usepackage{amsmath,amssymb,amsthm,mathtools}
\usepackage{enumitem}
\usepackage{booktabs}
\usepackage{graphicx}
\usepackage{microtype}
\usepackage[hidelinks]{hyperref}
\usepackage[a4paper,margin=1in]{geometry}

\newtheorem{theorem}{Theorem}[section]
\newtheorem{proposition}[theorem]{Proposition}
\newtheorem{lemma}[theorem]{Lemma}
\newtheorem{corollary}[theorem]{Corollary}
\newtheorem{remark}[theorem]{Remark}
\newtheorem{definition}[theorem]{Definition}
\newtheorem{example}[theorem]{Example}

\newcommand{\F}{\mathbb F}
\newcommand{\N}{\mathbb N}
\newcommand{\K}{\mathcal K}
\newcommand{\supp}{\operatorname{supp}}
\newcommand{\aco}{\operatorname{aco}}
\newcommand{\cl}{\operatorname{cl}}
\newcommand{\weak}{\sigma}
\newcommand{\EV}{\mathrm{EV}}
\newcommand{\ASCP}{\mathrm{ASCP}}
\newcommand{\wCSP}{\mathrm{wCSP}}
\newcommand{\Pin}{P_{\!\omega}}

\title{Pe{\l}czy\'nski $V^*$-Type Properties of Free Locally Convex Spaces:\\
A Complete Classification}
\author{Full solution packet, human-review draft}
\date{June 2026}

\begin{document}
\maketitle

\begin{abstract}
Let $L(X)$ be the free locally convex space over a Tychonoff space $X$, and let
$1\leq p\leq q\leq\infty$.  We give a complete answer to Problem~3.18 in
Gabriyelyan's paper on Gelfand--Phillips type properties\cite{Gabriyelyan2024GP}.  The properties
$V^*_{(p,q)}$, $\EV^*_{(p,q)}$, $sV^*_{(p,q)}$, and $s\EV^*_{(p,q)}$ are all
equivalent to the condition that every functionally bounded subset of $X$ is
finite.  In contrast, $L(X)$ has $wsV^*_{(p,q)}$ for every Tychonoff $X$, and it
has $ws\EV^*_{(p,q)}$ for every $X$ whenever $q<\infty$.  At the sole exceptional
endpoint $q=\infty$, the class of $p$-$(\EV^*)$ sets is exactly the class of
bounded sets.  Consequently, $ws\EV^*_p$ is equivalent to the weak Cauchy
subsequence property of $L(X)$.  We translate this last condition into an exact
sequential compactness condition for finitely supported probability measures on
compact closures of countable functionally bounded subsets of $X$ in the
Dieudonn\'e completion.
\end{abstract}

\section*{Source problem evidence}

\begin{center}
\includegraphics[width=\textwidth]{figures/open_problem_crop.png}
\end{center}

The source problem is Problem~3.18 in Gabriyelyan's paper
\cite{Gabriyelyan2024GP}.  It asks: for \(1\leq p\leq q\leq\infty\),
characterize the Tychonoff spaces \(X\) for which \(L(X)\) has one of the
\(V^\ast_{(p,q)}\)-type properties from Definition~2.5 of that paper.  The
definition includes \(V^\ast_{(p,q)}\), \(EV^\ast_{(p,q)}\),
\(sV^\ast_{(p,q)}\), \(sEV^\ast_{(p,q)}\), \(wsV^\ast_{(p,q)}\), and
\(wsEV^\ast_{(p,q)}\).

\section{Introduction and the classification theorem}

All locally convex spaces are Hausdorff and are over $\F=\mathbb R$ or
$\mathbb C$.  For a Tychonoff space $X$, every element of the free locally
convex space $L(X)$ has a unique representation
\[
  \chi=\sum_{j=1}^m a_jx_j,
\]
where the $x_j$ are distinct points of $X$.  We write
\[
  \|\chi\|_1=\sum_{j=1}^m |a_j|,
  \qquad
  \supp(\chi)=\{x_1,\dots,x_m\}.
\]
For $A\subseteq L(X)$, set
\[
  \supp(A)=\bigcup_{\chi\in A}\supp(\chi),
  \qquad
  C_A=\sup_{\chi\in A}\|\chi\|_1.
\]
The dual of $L(X)$ is $C(X)$, with the canonical pairing
\[
  \langle f,\chi\rangle=\sum_{j=1}^m a_jf(x_j).
\]

We recall the definitions needed for the statement.  A nonempty set $A$ in a
locally convex space $E$ is a $(p,q)$-$(V^*)$ set if, for every weakly
$p$-summable sequence $(u_n)$ in the strong dual $E'_\beta$, the sequence
\[
  \left(\sup_{a\in A}|\langle u_n,a\rangle|\right)_n
\]
belongs to $\ell_q$ when $q<\infty$, and belongs to $c_0$ when $q=\infty$.
A $(p,q)$-$(\EV^*)$ set is defined in the same way, but only equicontinuous
weakly $p$-summable sequences are tested.  The properties $V^*_{(p,q)}$ and
$\EV^*_{(p,q)}$ require these sets to be relatively weakly compact; the
properties $sV^*_{(p,q)}$ and $s\EV^*_{(p,q)}$ require relative weak sequential
compactness; and $wsV^*_{(p,q)}$ and $ws\EV^*_{(p,q)}$ require weak sequential
precompactness, meaning that every sequence has a weakly Cauchy subsequence.
When $q=\infty$, the second subscript is omitted.

The following theorem is the full classification proved below.

\begin{theorem}[Complete classification]\label{thm:classification}
Let $X$ be a Tychonoff space and let $1\leq p\leq q\leq\infty$.
\begin{enumerate}[label=\textup{(\Roman*)}]
\item\label{cl:strong}
The following assertions are equivalent:
\begin{enumerate}[label=\textup{(\alph*)}]
\item every functionally bounded subset of $X$ is finite;
\item $L(X)$ has $V^*_{(p,q)}$;
\item $L(X)$ has $\EV^*_{(p,q)}$;
\item $L(X)$ has $sV^*_{(p,q)}$;
\item $L(X)$ has $s\EV^*_{(p,q)}$.
\end{enumerate}
\item\label{cl:wsv}
The space $L(X)$ has $wsV^*_{(p,q)}$ for every Tychonoff space $X$.
\item\label{cl:wsevfinite}
If $q<\infty$, then $L(X)$ has $ws\EV^*_{(p,q)}$ for every Tychonoff space $X$.
\item\label{cl:endpoint}
For $q=\infty$, the following assertions are equivalent, and do not depend on
$p$:
\begin{enumerate}[label=\textup{(\alph*)}]
\item $L(X)$ has $ws\EV^*_p$;
\item every bounded sequence in $L(X)$ has a weakly Cauchy subsequence;
\item $X$ has the atomic sequential compactness property $\ASCP$ of
Definition~\ref{def:ascp} below.
\end{enumerate}
\end{enumerate}
\end{theorem}

Thus only the equicontinuous endpoint $ws\EV^*_{(p,\infty)}$ depends on $X$ in
a nontrivial way.  The condition in Theorem~\ref{thm:classification}\ref{cl:endpoint}
will be expressed solely in terms of probability measures over compact closures
of countable functionally bounded subsets of $X$.

\section*{Proof intuition}

The four compactness and weak-sequential-compactness properties are governed by
one simple obstruction: an infinite functionally bounded set in \(X\) supports
weighted finite partial sums in \(L(X)\) which are \(V^\ast\)-small but converge
pointwise to an infinite measure outside \(L(X)\).  This proves the necessity of
``functionally bounded sets are finite'', while the converse follows because
bounded \(EV^\ast\)-sets then live in finite-dimensional coordinate spans.

The weak sequentially precompact variants require a different mechanism.  If a
bounded subset of \(L(X)\) contains no weakly Cauchy subsequence, its supports
sit in a compact \(K\subseteq\mu X\), and the corresponding measures fail weak
compactness in \(M(K)\).  The Dieudonne--Grothendieck disjoint-open-set
criterion gives disjoint tests on \(K\).  A telescoping extension turns these
tests into a weakly \(1\)-summable sequence in \(C(X)_\beta\) which detects the
failure uniformly.  This immediately forces every \(V^\ast\)-set to be weakly
sequentially precompact and, after a \(c_0\)-damping trick, gives the
equicontinuous result for finite \(q\).  The remaining endpoint
\(wsEV^\ast_p\) is exactly the assertion that bounded sequences in \(L(X)\) have
weakly Cauchy subsequences; this is reformulated as a sequential compactness
condition for finitely supported probability measures on compact closures of
countable functionally bounded subsets of \(X\).

\section{Free locally convex preliminaries}

Let $\mu X$ denote the Dieudonn\'e completion of $X$.  We use the following
standard bounded-set criterion for free locally convex spaces
\cite[Proposition~2.7]{Gabriyelyan2019}.

\begin{proposition}[Bounded sets in $L(X)$]\label{prop:bounded}
For a nonempty set $B\subseteq L(X)$, the following are equivalent:
\begin{enumerate}[label=\textup{(\roman*)}]
\item $B$ is bounded in $L(X)$;
\item $\supp(B)$ has compact closure in $\mu X$ and $C_B<\infty$;
\item $\supp(B)$ is functionally bounded in $X$ and $C_B<\infty$.
\end{enumerate}
\end{proposition}

We record several consequences in a form used repeatedly below.

\begin{lemma}[Restriction to compact closures]\label{lem:restriction}
Let $D\subseteq X$ be functionally bounded and put
\[
  K_D=\cl_{\mu X}D.
\]
Then $K_D$ is compact.  Every $f\in C(X)$ has a unique continuous extension
$\bar f\in C(\mu X)$, and the restriction operator
\[
  R_D:C(X)\longrightarrow C(K_D),
  \qquad R_Df=\bar f|_{K_D},
\]
is onto.  Moreover,
\[
  \|R_Df\|_\infty=\sup_{x\in D}|f(x)|.
\]
\end{lemma}

\begin{proof}
A subset of $X$ is functionally bounded if and only if its closure in $\mu X$
is compact; this is part of the standard theory of the universal uniformity and
is also used in Proposition~\ref{prop:bounded}.  Every continuous scalar-valued
function on $X$ is uniformly continuous for the universal uniformity, hence
extends uniquely to $\mu X$.

It remains to prove surjectivity.  A compact subset of a Tychonoff space is
$C$-embedded.  Indeed, regard $K_D$ as a closed subset of the Stone--\v Cech
compactification $\beta(\mu X)$, extend a scalar-valued continuous function on
$K_D$ by the Tietze extension theorem, and restrict the extension to $\mu X$
and then to $X$.  Finally, $D$ is dense in $K_D$, so the displayed norm identity
holds.
\end{proof}

\begin{lemma}[The strong dual of $L(X)$]\label{lem:strongdual}
The strong topology $\beta(C(X),L(X))$ is generated by the seminorms
\[
  p_D(f)=\sup_{x\in D}|f(x)|,
\]
where $D$ ranges over the functionally bounded subsets of $X$.
Furthermore, for every
\(
  \Lambda\in(C(X)_\beta)'
\)
there are a functionally bounded set $D\subseteq X$ and a finite regular Borel
measure $\nu\in M(K_D)=C(K_D)'$ such that
\[
  \Lambda(f)=\int_{K_D}R_Df\,d\nu
  \qquad(f\in C(X)).
\]
\end{lemma}

\begin{proof}
If $D$ is functionally bounded, then $\{x:x\in D\}$ is bounded in $L(X)$ by
Proposition~\ref{prop:bounded}; hence $p_D$ is a strong-dual seminorm.
Conversely, if $B\subseteq L(X)$ is bounded, then
\[
  \sup_{\chi\in B}|\langle f,\chi\rangle|
  \leq C_B\sup_{x\in\supp(B)}|f(x)|,
\]
and $\supp(B)$ is functionally bounded.  This proves the first assertion.

If $\Lambda$ is continuous, there are a functionally bounded $D$ and $c>0$
such that
\[
  |\Lambda(f)|\leq c p_D(f)
  \qquad(f\in C(X));
\]
a finite union absorbs the finitely many seminorms arising from a basic
neighborhood.  Thus $\Lambda$ vanishes on $\ker R_D$.  By
Lemma~\ref{lem:restriction}, it induces a bounded linear functional on
$C(K_D)$.  The Riesz representation theorem gives the required measure $\nu$.
\end{proof}

We shall also use two standard Banach-space facts.  First, by the
Eberlein--\v Smulian theorem, a relatively weakly compact subset of a Banach
space is relatively weakly sequentially compact.  Second, we use the following
form of the Dieudonn\'e--Grothendieck criterion
\cite{Wells1969}.

\begin{theorem}[Dieudonn\'e--Grothendieck criterion]\label{thm:DG}
Let $K$ be compact and let $H\subseteq M(K)$ be norm bounded.  Then $H$ is
relatively weakly compact if and only if, for every sequence $(U_n)$ of pairwise
disjoint open subsets of $K$,
\[
  \sup_{\mu\in H}|\mu(U_n)|\longrightarrow0.
\]
\end{theorem}

\section{The four compactness and sequential-compactness variants}

We begin with the universal obstruction supplied by an infinite functionally
bounded set.

\begin{lemma}[Weighted partial sums]\label{lem:weighted-v}
Let $D=\{x_i:i\in\N\}$ be an infinite functionally bounded subset of $X$, with
all $x_i$ distinct, and let $w_i=2^{-i-1}$.  Put
\[
  a_m=\sum_{i=0}^m w_i x_i,
  \qquad
  A=\{a_m:m\in\N\}.
\]
Then $A$ is a $(p,q)$-$(V^*)$ set in $L(X)$ for every
$1\leq p\leq q\leq\infty$.
\end{lemma}

\begin{proof}
The restriction operator
\[
  R:C(X)_\beta\longrightarrow\ell_\infty,
  \qquad Rf=(f(x_i))_{i\in\N},
\]
is continuous by Lemma~\ref{lem:strongdual}.  Define
\[
  H:\ell_\infty\longrightarrow\ell_\infty,
  \qquad
  H(y)(m)=\sum_{i=0}^m w_i y(i).
\]

Assume first that $1\leq p<\infty$ and let $(y_k)$ be weakly $p$-summable in
$\ell_\infty$.  The operator
\[
  S:(\ell_\infty)'\longrightarrow\ell_p,
  \qquad S(\varphi)=(\varphi(y_k))_k,
\]
is well defined and has closed graph, hence is bounded.  If $e_i'$ is the
$i$th coordinate functional, then
\[
  \big\|(y_k(i))_k\big\|_{\ell_p}=\|S(e_i')\|_{\ell_p}\leq\|S\|.
\]
Minkowski's inequality therefore gives
\[
\begin{split}
  \left(\sum_k\|H(y_k)\|_\infty^p\right)^{1/p}
  &\leq
  \sum_i w_i\left(\sum_k|y_k(i)|^p\right)^{1/p}\\
  &\leq \|S\|\sum_iw_i<\infty.
\end{split}
\]
Thus $(\|H(y_k)\|_\infty)_k\in\ell_p$, and hence it belongs to $\ell_q$ when
$q<\infty$ and to $c_0$ when $q=\infty$.

For $p=\infty$, the finite-rank operators
\[
  H_N(y)(m)=\sum_{i=0}^{\min\{m,N\}}w_i y(i)
\]
converge to $H$ in operator norm.  Hence $H$ is compact and sends weakly null
sequences to norm-null sequences.

Now let $(f_k)$ be weakly $p$-summable in $C(X)_\beta$.  Then $(Rf_k)$ is
weakly $p$-summable in $\ell_\infty$, and
\[
  \sup_{a\in A}|\langle f_k,a\rangle|
  =\sup_m\left|\sum_{i=0}^m w_i f_k(x_i)\right|
  =\|H(Rf_k)\|_\infty.
\]
The preceding estimates prove the assertion.
\end{proof}

\begin{lemma}[The weighted partial sums do not compactify in $L(X)$]
\label{lem:weighted-noncompact}
With the notation of Lemma~\ref{lem:weighted-v}, the set $A$ is neither
relatively weakly compact nor relatively weakly sequentially compact in $L(X)$.
\end{lemma}

\begin{proof}
Define
\[
  \mu(f)=\sum_{i=0}^\infty w_i f(x_i),
  \qquad f\in C(X).
\]
The series converges absolutely because $D$ is functionally bounded.  It also
defines an element of $(C(X)_\beta)'$, since
\[
  |\mu(f)|\leq\left(\sum_iw_i\right)p_D(f).
\]
For every $f\in C(X)$,
\[
  \langle f,a_m\rangle\longrightarrow\mu(f).
\]

The functional $\mu$ is not represented by an element of $L(X)$.  Otherwise,
write $\mu=\sum_{j=1}^r c_jy_j$.  Choose $N$ with
$x_N\notin\{y_1,\dots,y_r\}$.  Complete regularity supplies
$f\in C(X,[0,1])$ such that $f(x_N)=1$ and $f$ vanishes on
\[
  \{y_1,\dots,y_r,x_0,\dots,x_{N-1}\}.
\]
The finite representation gives $\mu(f)=0$, whereas the series gives
$\mu(f)\geq w_N>0$, a contradiction.

Let
\[
  J:L(X)_w\longrightarrow\F^{C(X)},
  \qquad J(\chi)=(\langle f,\chi\rangle)_{f\in C(X)},
\]
be the canonical weak-topology embedding.  If the weak closure of $A$ were
compact, its image under $J$ would be compact and therefore closed in the
Hausdorff product $\F^{C(X)}$.  Since $J(a_m)$ converges coordinatewise to
$(\mu(f))_{f\in C(X)}$, that point would have to belong to $J(L(X))$, contrary
to the preceding paragraph.  Thus $A$ is not relatively weakly compact.
Every subsequence of $(a_m)$ has the same coordinatewise limit $\mu$, so no
subsequence converges weakly in $L(X)$.  Hence $A$ is not relatively weakly
sequentially compact.
\end{proof}

We can now settle the first four properties.

\begin{theorem}\label{thm:four}
Let $1\leq p\leq q\leq\infty$.  The following are equivalent:
\begin{enumerate}[label=\textup{(\roman*)}]
\item every functionally bounded subset of $X$ is finite;
\item $L(X)$ has $V^*_{(p,q)}$;
\item $L(X)$ has $\EV^*_{(p,q)}$;
\item $L(X)$ has $sV^*_{(p,q)}$;
\item $L(X)$ has $s\EV^*_{(p,q)}$.
\end{enumerate}
\end{theorem}

\begin{proof}
Assume first that every functionally bounded subset of $X$ is finite.  Every
$(p,q)$-$(V^*)$ set is a $(p,q)$-$(\EV^*)$ set, and every
$(p,q)$-$(\EV^*)$ set is bounded
\cite[Lemma~7.2]{GabriyelyanPelczynski}.  If $B$ is such a set, then
Proposition~\ref{prop:bounded} gives $C_B<\infty$ and says that
$F=\supp(B)$ is functionally bounded.  Thus $F$ is finite and
\[
  B\subseteq
  \left\{\sum_{x\in F}a_xx:\sum_{x\in F}|a_x|\leq C_B\right\}.
\]
The set on the right is compact in the finite-dimensional space
$\operatorname{span}(F)$.  Hence $B$ is relatively compact, relatively weakly
compact, relatively weakly sequentially compact, and weakly sequentially
precompact.  This proves all positive implications.

Conversely, if $X$ contains an infinite functionally bounded subset,
Lemmas~\ref{lem:weighted-v} and~\ref{lem:weighted-noncompact} provide a
$(p,q)$-$(V^*)$ set which is neither relatively weakly compact nor relatively
weakly sequentially compact.  Since every $(V^*)$ set is an $(\EV^*)$ set, the
same witness defeats all four properties.
\end{proof}

\section{A detector for failure of weak sequential precompactness}

The new ingredient needed for the weak variants is the following lemma.  It
turns a bounded sequence without a weakly Cauchy subsequence into a weakly
$1$-summable sequence of continuous functions that detects that failure
uniformly.

\begin{lemma}[Disjointization and telescoping detector]\label{lem:detector}
Let $A\subseteq L(X)$ be bounded and suppose that $A$ contains a sequence with
no weakly Cauchy subsequence.  Then there are $\eta>0$ and functions
$h_n\in C(X,[0,1])$ such that
\begin{enumerate}[label=\textup{(\roman*)}]
\item $(h_n)$ is weakly $1$-summable in $C(X)_\beta$;
\item $\displaystyle \sum_{n=0}^\infty h_n(x)\leq1$ for every $x\in X$;
\item $\displaystyle \sup_{\chi\in A}|\langle h_n,\chi\rangle|\geq\eta$ for
every $n$.
\end{enumerate}
\end{lemma}

\begin{proof}
Choose a sequence $(\chi_m)\subseteq A$ with no weakly Cauchy subsequence and
put
\[
  S=\bigcup_m\supp(\chi_m),
  \qquad K=\cl_{\mu X}S.
\]
The set $S$ is countable and functionally bounded, and the total variations
$\|\chi_m\|_1$ are uniformly bounded.  We therefore regard the $\chi_m$ as a
bounded sequence in $M(K)=C(K)'$.

The set $H=\{\chi_m:m\in\N\}$ is not relatively weakly compact in $M(K)$.
Indeed, otherwise Eberlein--\v Smulian would give a subsequence converging weakly
in $M(K)$, hence converging pointwise on $C(K)$.  By
Lemma~\ref{lem:restriction}, every member of $C(K)$ is the restriction of a
member of $C(X)$, so the same subsequence would be weakly Cauchy in $L(X)$, a
contradiction.

Apply Theorem~\ref{thm:DG}.  There are pairwise disjoint open sets
$U_n\subseteq K$, measures $\mu_n\in H$, and $\varepsilon>0$ such that
\[
  |\mu_n(U_n)|>3\varepsilon
  \qquad(n\in\N).
\]
By regularity, choose a compact set $F_n\subseteq U_n$ such that
$|\mu_n|(U_n\setminus F_n)<\varepsilon$.  Normality of $K$ gives
$g_n\in C(K,[0,1])$ with $g_n=1$ on $F_n$ and
$\supp(g_n)\subseteq U_n$.  Then
\[
  |\mu_n(g_n)|>2\varepsilon.
\]
The supports of the $g_n$ are pairwise disjoint.

Since $K$ is compact in the Tychonoff space $\mu X$, it is $C$-embedded.
Extend each $g_n$ to a function $u_n\in C(\mu X,[0,1])$.  Define
\[
  \widetilde h_0=u_0,
  \qquad
  \widetilde h_n=u_n\prod_{j<n}(1-u_j)
  \quad(n\geq1),
\]
and let $h_n=\widetilde h_n|_X$.  Because the supports of the $g_n$ are
disjoint,
\[
  \widetilde h_n|_K=g_n.
\]
Moreover, for every $x\in\mu X$ and every $N$,
\[
  \sum_{n=0}^N\widetilde h_n(x)
  =1-\prod_{n=0}^N(1-u_n(x))\leq1.
\]
This proves (ii), and since each $\mu_n$ is supported on $K$,
\[
  \sup_{\chi\in A}|\langle h_n,\chi\rangle|
  \geq|\langle h_n,\mu_n\rangle|
  =|\mu_n(g_n)|>2\varepsilon.
\]
Thus (iii) holds with $\eta=2\varepsilon$.

It remains to prove (i).  Let $\Lambda\in(C(X)_\beta)'$.  By
Lemma~\ref{lem:strongdual}, there are a functionally bounded set $D$ and
$\nu\in M(K_D)$ such that
\[
  \Lambda(f)=\int_{K_D}\bar f\,d\nu.
\]
Since the same telescoping estimate holds on all of $\mu X$, Tonelli's theorem
gives
\[
\begin{split}
  \sum_n|\Lambda(h_n)|
  &\leq\sum_n\int_{K_D}\widetilde h_n\,d|\nu|\\
  &=\int_{K_D}\sum_n\widetilde h_n\,d|\nu|
  \leq\|\nu\|.
\end{split}
\]
Thus $(\Lambda(h_n))\in\ell_1$ for every
$\Lambda\in(C(X)_\beta)'$, which is precisely weak $1$-summability.
\end{proof}

\section{The universal weak variants}

\begin{theorem}[Universality of $wsV^*$]\label{thm:wsv}
For every Tychonoff space $X$ and every $1\leq p\leq q\leq\infty$, the space
$L(X)$ has $wsV^*_{(p,q)}$.
\end{theorem}

\begin{proof}
Let $A$ be a $(p,q)$-$(V^*)$ set.  It is bounded because it is also a
$(p,q)$-$(\EV^*)$ set and every such set is bounded.  Suppose that $A$ is not
weakly sequentially precompact.  Apply Lemma~\ref{lem:detector}.  The resulting
sequence $(h_n)$ is weakly $1$-summable, hence weakly $p$-summable, in
$C(X)_\beta$.  However,
\[
  \sup_{\chi\in A}|\langle h_n,\chi\rangle|\geq\eta
  \qquad(n\in\N),
\]
so this sequence of suprema belongs neither to $\ell_q$ when $q<\infty$ nor to
$c_0$ when $q=\infty$.  This contradicts the $(p,q)$-$(V^*)$ property of $A$.
\end{proof}

\begin{theorem}[Universality of finite-$q$ $ws\EV^*$]\label{thm:wsevfinite}
Let $1\leq p\leq q<\infty$.  Then $L(X)$ has $ws\EV^*_{(p,q)}$ for every
Tychonoff space $X$.
\end{theorem}

\begin{proof}
Let $A$ be a $(p,q)$-$(\EV^*)$ set and suppose that it is not weakly
sequentially precompact.  It is bounded, so Lemma~\ref{lem:detector} provides
$\eta>0$ and a weakly $1$-summable sequence $(h_n)$ satisfying
\[
  \sum_nh_n(x)\leq1
  \quad(x\in X),
  \qquad
  \sup_{\chi\in A}|\langle h_n,\chi\rangle|\geq\eta.
\]
Choose positive scalars $b_n\to0$ with $(b_n)\notin\ell_q$, for instance
$b_n=(n+1)^{-1/q}$, and put $v_n=b_nh_n$.  The sequence $(v_n)$ remains weakly
$1$-summable, hence weakly $p$-summable.

We claim that $\{v_n:n\in\N\}$ is equicontinuous in $L(X)'$.  The map
\[
  \Phi:X\longrightarrow c_0,
  \qquad
  \Phi(x)=(v_n(x))_n,
\]
is continuous: each coordinate is continuous, while
\[
  \sup_{x\in X}|v_n(x)|\leq b_n\longrightarrow0,
\]
so the tails are uniformly small.  By the universal property of $L(X)$,
$\Phi$ extends to a continuous linear operator $T:L(X)\to c_0$.  Since
$v_n=T'(e_n')$ and the coordinate functionals $(e_n')$ form an equicontinuous
subset of $c_0'$, the sequence $(v_n)$ is equicontinuous.

On the other hand,
\[
  \sup_{\chi\in A}|\langle v_n,\chi\rangle|
  \geq\eta b_n,
\]
and $(b_n)\notin\ell_q$.  This contradicts the $(p,q)$-$(\EV^*)$ property of
$A$.
\end{proof}

The damping $b_n\to0$ is exactly what creates equicontinuity.  It also explains
why the preceding argument cannot contradict an $(p,\infty)$-$(\EV^*)$
condition: the lower bound $\eta b_n$ necessarily tends to zero.  The endpoint
$q=\infty$ therefore requires a separate analysis.

\section{The exceptional endpoint \texorpdfstring{$ws\EV^*_p$}{wsEV-p}}

\begin{proposition}[At $q=\infty$, $p$-$(\EV^*)$ means bounded]
\label{prop:evbounded}
For every Tychonoff space $X$ and every $p\in[1,\infty]$, a subset of $L(X)$ is
a $p$-$(\EV^*)$ set if and only if it is bounded.
\end{proposition}

\begin{proof}
Every $p$-$(\EV^*)$ set is bounded
\cite[Lemma~7.2]{GabriyelyanPelczynski}.  Conversely, let $B\subseteq L(X)$ be
bounded.  Put
\[
  D=\supp(B),
  \qquad C=C_B<\infty.
\]
Let $(f_n)$ be an equicontinuous weakly $p$-summable sequence in
$C(X)_\beta$.  Weak $p$-summability implies
\[
  f_n(x)\longrightarrow0
  \qquad(x\in X),
\]
because evaluation at $x$ is a continuous functional on $C(X)_\beta$.
Equicontinuity says that
\[
  \chi\longmapsto\sup_n|\langle f_n,\chi\rangle|
\]
is a continuous seminorm on $L(X)$.  Hence
\[
  T:L(X)\longrightarrow\ell_\infty,
  \qquad T(\chi)=(\langle f_n,\chi\rangle)_n,
\]
is continuous.  Its restriction to $X$ takes values in $c_0$, so
\[
  \Phi:X\longrightarrow c_0,
  \qquad \Phi(x)=(f_n(x))_n,
\]
is continuous.

The set $D$ is functionally bounded.  Therefore $\Phi(D)$ is functionally
bounded in the metric space $c_0$, and hence has compact closure $Q$ in $c_0$.
Compact subsets of $c_0$ have uniformly vanishing coordinates, so
\[
  \sup_{x\in D}|f_n(x)|
  =\sup_{z\in\Phi(D)}|z_n|
  \leq\sup_{z\in Q}|z_n|
  \longrightarrow0.
\]
Consequently,
\[
  \sup_{\chi\in B}|\langle f_n,\chi\rangle|
  \leq C\sup_{x\in D}|f_n(x)|
  \longrightarrow0.
\]
Thus $B$ is a $p$-$(\EV^*)$ set.
\end{proof}

\begin{corollary}\label{cor:endpoint-wcsp}
For every $p\in[1,\infty]$,
\[
  L(X)\text{ has }ws\EV^*_p
  \quad\Longleftrightarrow\quad
  \text{every bounded sequence in }L(X)\text{ has a weakly Cauchy subsequence}.
\]
In particular, the endpoint property is independent of $p$.
\end{corollary}

\begin{proof}
By Proposition~\ref{prop:evbounded}, the $p$-$(\EV^*)$ subsets of $L(X)$ are
exactly its bounded subsets.
\end{proof}

We next translate Corollary~\ref{cor:endpoint-wcsp} into a condition on $X$.
For a countable functionally bounded set $S\subseteq X$, put
\[
  K_S=\cl_{\mu X}S
\]
and let
\[
  \Pin(S)=
  \left\{
    \sum_{j=1}^m t_j\delta_{s_j}:
    s_j\in S,\ t_j\geq0,\ \sum_{j=1}^m t_j=1
  \right\}
  \subseteq P(K_S).
\]
Here $P(K_S)$ is the weak-star compact space of Radon probability measures on
$K_S$.

\begin{definition}[Atomic sequential compactness property]\label{def:ascp}
A Tychonoff space $X$ has property $\ASCP$ if, for every countable
functionally bounded set $S\subseteq X$, every sequence in $\Pin(S)$ has a
weak-star convergent subsequence in $P(K_S)$.
\end{definition}

\begin{lemma}[Probability measures suffice]\label{lem:probabilities}
Fix a countable functionally bounded $S\subseteq X$.  The following are
equivalent:
\begin{enumerate}[label=\textup{(\roman*)}]
\item every sequence in $\Pin(S)$ has a weak-star convergent subsequence in
$P(K_S)$;
\item every norm-bounded sequence of finitely supported real signed measures
carried by $S$ has a weak-star convergent subsequence in $M(K_S)$;
\item every norm-bounded sequence of finitely supported complex measures
carried by $S$ has a weak-star convergent subsequence in $M(K_S)$.
\end{enumerate}
\end{lemma}

\begin{proof}
Only \(\textup{(i)}\Rightarrow\textup{(ii)}\) requires an argument.  Let $(\mu_n)$ be uniformly
variation bounded and write the Jordan decomposition
\[
  \mu_n=\mu_n^+-\mu_n^-.
\]
After passing to a subsequence, the masses
$\alpha_n=\mu_n^+(K_S)$ and $\beta_n=\mu_n^-(K_S)$ converge.  If the limit of
$\alpha_n$ is zero, then $\|\mu_n^+\|\to0$; otherwise, normalize
$\mu_n^+/\alpha_n\in\Pin(S)$ and use (i).  Treat the negative parts in the same
way and take a diagonal subsequence.  This proves (ii).  Splitting a complex
measure into its real and imaginary parts gives (iii).  The reverse implications
are immediate.
\end{proof}

\begin{theorem}[Topological-measure characterization of the endpoint]
\label{thm:ascp}
For a Tychonoff space $X$, the following are equivalent:
\begin{enumerate}[label=\textup{(\roman*)}]
\item $L(X)$ has $ws\EV^*_p$ for one, equivalently every, $p\in[1,\infty]$;
\item every bounded sequence in $L(X)$ has a weakly Cauchy subsequence;
\item $X$ has $\ASCP$.
\end{enumerate}
\end{theorem}

\begin{proof}
The equivalence \(\textup{(i)}\Leftrightarrow\textup{(ii)}\) is
Corollary~\ref{cor:endpoint-wcsp}.

Assume (ii), let $S$ be countable and functionally bounded, and let
$(\mu_n)\subseteq\Pin(S)$.  By Proposition~\ref{prop:bounded}, $(\mu_n)$ is a
bounded sequence in $L(X)$.  It therefore has a weakly Cauchy subsequence.
Every $g\in C(K_S)$ extends to a member of $C(X)$ by
Lemma~\ref{lem:restriction}; hence this subsequence is weak-star Cauchy on
$C(K_S)$.  Its pointwise limit is a positive norm-one functional on $C(K_S)$,
so it is a member of $P(K_S)$.  Thus $X$ has $\ASCP$.

Conversely, assume $\ASCP$ and let $(\chi_n)$ be a bounded sequence in $L(X)$.
Set
\[
  S=\bigcup_n\supp(\chi_n).
\]
Then $S$ is countable and functionally bounded, and the variations
$\|\chi_n\|_1$ are uniformly bounded.  By Lemma~\ref{lem:probabilities}, a
subsequence converges weak-star in $M(K_S)$.  Therefore its values converge on
every restriction $R_Sf$, $f\in C(X)$, so the same subsequence is weakly Cauchy
in $L(X)$.  This proves (ii).
\end{proof}

The restriction-space formulation from the earlier partial result becomes
especially transparent here.

\begin{proposition}[Restriction-space form]\label{prop:restriction-space}
For every countable functionally bounded $S\subseteq X$, the space
\[
  E_S=\overline{\{f|_S:f\in C(X)\}}^{\|\cdot\|_\infty}
  \subseteq\ell_\infty(S)
\]
is canonically isometric to $C(K_S)$.  In fact the displayed restriction range
is already closed.  Property $\ASCP$ says exactly that
\[
  \operatorname{co}\{\delta_s:s\in S\}\subseteq E_S'
\]
is relatively weak-star sequentially compact for every such $S$.
\end{proposition}

\begin{proof}
By Lemma~\ref{lem:restriction}, restriction from $C(X)$ to $S$ has exactly the
same range as restriction from $C(K_S)$ to its dense subset $S$.  The latter
map is an isometry, so its range is complete and therefore closed.
\end{proof}

\section{Consequences and examples}

\begin{corollary}\label{cor:sufficient}
The endpoint $ws\EV^*_p$ holds, and hence all weak variants hold, under any of
the following conditions:
\begin{enumerate}[label=\textup{(\roman*)}]
\item for every countable functionally bounded $S\subseteq X$, the compactum
$K_S$ is metrizable;
\item for every such $S$, the compactum $K_S$ is Rosenthal compact;
\item $X$ is metrizable.
\end{enumerate}
\end{corollary}

\begin{proof}
If $K_S$ is metrizable, then $P(K_S)$ is compact metrizable and therefore
sequentially compact.  If $K_S$ is Rosenthal compact, then $P(K_S)$ is
Rosenthal compact by Godefroy's theorem \cite{Godefroy1980}; Rosenthal compacta
are Fr\'echet--Urysohn by the Bourgain--Fremlin--Talagrand theorem
\cite{BFT1978}, and compact Fr\'echet--Urysohn spaces are sequentially compact.
In either case $\ASCP$ follows.

Every metrizable space is a $\mu$-space: functionally bounded subsets have
compact closure.  Hence, for countable functionally bounded $S$ in a metrizable
$X$, the compactum $K_S=\cl_XS$ is metrizable.
\end{proof}

\begin{example}[Separation from the four stronger variants]
If $X=[0,1]$, then $L(X)$ has every $wsV^*_{(p,q)}$ and every
$ws\EV^*_{(p,q)}$, but it has none of
\[
  V^*_{(p,q)},\quad \EV^*_{(p,q)},\quad
  sV^*_{(p,q)},\quad s\EV^*_{(p,q)}.
\]
Indeed, the weak properties follow from Corollary~\ref{cor:sufficient}, while
the four stronger properties fail by Theorem~\ref{thm:four} because $[0,1]$ is
an infinite functionally bounded set.
\end{example}

\begin{example}[The endpoint is not universal]\label{ex:betan}
Let $X=\beta\N$.  The sequence $(\delta_n)$ is bounded in $L(\beta\N)$, but it
has no weakly Cauchy subsequence.  Indeed, given a subsequence $(n_k)$, choose
$a\in\ell_\infty=C(\beta\N)$ with $a(n_k)=(-1)^k$.  Then
$\langle a,\delta_{n_k}\rangle=(-1)^k$ is not Cauchy.  Thus $\beta\N$ fails
$\ASCP$, and
\[
  L(\beta\N)\text{ fails }ws\EV^*_p
  \quad\text{for every }p.
\]
Nevertheless, Theorems~\ref{thm:wsv} and~\ref{thm:wsevfinite} show that
$L(\beta\N)$ has all $wsV^*_{(p,q)}$ properties and all
$ws\EV^*_{(p,q)}$ properties with $q<\infty$.
\end{example}

\begin{remark}[Why the endpoint does not collapse to a standard adjective]
The condition $\ASCP$ is an exact condition on $X$, but it need not reduce to
metrizability of the compacta $K_S$.  Sequential properties of spaces of
probability measures over nonmetrizable compacta are delicate.  For example,
it remains open in general whether $P(K)$ must be sequentially compact whenever
$K$ is sequentially compact; see \cite[Problem~7.8(4)]{Plebanek2024}.  Thus an
attempt to replace $\ASCP$ by a familiar one-word compactness condition would
encounter a genuinely nontrivial problem in the topology of $P(K)$ spaces.
\end{remark}

\section{Conclusion}

Theorem~\ref{thm:classification} follows by combining
Theorems~\ref{thm:four}, \ref{thm:wsv}, \ref{thm:wsevfinite}, and
\ref{thm:ascp}.  In compact form, the answer is
\begin{center}
\begin{tabular}{@{}ll@{}}
\toprule
Property of $L(X)$ & Exact condition on $X$\\
\midrule
$V^*_{(p,q)},\ \EV^*_{(p,q)},\ sV^*_{(p,q)},\ s\EV^*_{(p,q)}$
& Every functionally bounded subset of $X$ is finite\\
$wsV^*_{(p,q)}$ & Always\\
$ws\EV^*_{(p,q)}$, $q<\infty$ & Always\\
$ws\EV^*_p$ & $\ASCP$\\
\bottomrule
\end{tabular}
\end{center}
The central new mechanism is Lemma~\ref{lem:detector}: the
Dieudonn\'e--Grothendieck disjoint-open-set criterion is combined with a
telescoping extension to turn failure of weak sequential precompactness into a
weakly $1$-summable detector in the strong dual of $L(X)$.

\paragraph{Verification status.}
The weighted-partial-sum argument for the four stronger variants was already
present in the supplied partial packet.  The detector lemma and its consequences
for the weak variants are the additional part of this proof draft.  The argument
has been written so that all nonstandard steps are isolated in
Lemmas~\ref{lem:restriction}, \ref{lem:strongdual}, and~\ref{lem:detector}.
The local run registry was checked for prior packets on arXiv:2406.07178, and a
targeted web/arXiv search on June~18, 2026 for Problem~3.18, arXiv:2406.07178,
and the phrases ``\(V^\ast\)'', ``\(wsEV^\ast\)'', and ``free locally convex''
located the source problem and the background Pe{\l}czy\'nski \(V^\ast\) paper
but no separate published solution of Problem~3.18.  This is a bounded novelty
check, not a formal exhaustive review.  Human review should pay particular
attention to the detector lemma, the use of the Dieudonne--Grothendieck
criterion, and the endpoint identification of \(p\)-\((EV^\ast)\) sets with
bounded sets.

\begin{thebibliography}{99}

\bibitem{BFT1978}
J.~Bourgain, D.~H. Fremlin, and M.~Talagrand,
\newblock Pointwise compact sets of measurable functions,
\newblock \emph{Amer. J. Math.} \textbf{100} (1978), 845--886.

\bibitem{BanakhGabriyelyan2022}
T.~Banakh and S.~Gabriyelyan,
\newblock On free locally convex spaces,
\newblock \emph{Filomat} \textbf{36} (2022), no.~18, 6393--6401.

\bibitem{Godefroy1980}
G.~Godefroy,
\newblock Compacts de Rosenthal,
\newblock \emph{Pacific J. Math.} \textbf{91} (1980), 293--306.

\bibitem{Gabriyelyan2019}
S.~Gabriyelyan,
\newblock Locally convex properties of free locally convex spaces,
\newblock \emph{J. Math. Anal. Appl.} \textbf{480} (2019), 123453.

\bibitem{Gabriyelyan2024GP}
S.~Gabriyelyan,
\newblock Gelfand--Phillips type properties of locally convex spaces,
\newblock \emph{Mathematics} \textbf{12} (2024), no.~22, 3537.

\bibitem{GabriyelyanPelczynski}
S.~Gabriyelyan,
\newblock Pe{\l}czy\'nski's type sets and Pe{\l}czy\'nski's geometrical properties
of locally convex spaces,
\newblock arXiv:2402.08860 (2024).

\bibitem{Plebanek2024}
G.~Plebanek,
\newblock A survey on topological properties of $P(K)$ spaces,
\newblock \emph{Collect. Math.} (2024),
\newblock \href{https://doi.org/10.1007/s11537-024-2410-y}{doi:10.1007/s11537-024-2410-y}.

\bibitem{Wells1969}
B.~B. Wells, Jr.,
\newblock Weak compactness of measures,
\newblock \emph{Proc. Amer. Math. Soc.} \textbf{20} (1969), 124--130.

\end{thebibliography}

\end{document}
